% Sec. VII Related Work, Sec. VIII Limitations, Sec. IX Conclusion, Artifact Availability.
\section{Related Work}\label{sec:related}

Grundmann and Hartenstein~\cite{grundmann2026} model-check a stepwise refinement of the
Lightning Network in TLA+; model checking explores a bounded state space. Our \tamarin{} proofs
are symbolic and unbounded: subject to the restrictions of our model
(Sec.~\ref{sec:restrictions}), including three liveness restrictions, they cover any number of
parties, channels, channel updates and payments over routes of any length, and bypass resistance
holds at every position of a route for single-offer payments (Sec.~\ref{sec:bypass}).
\cite{litos2026} presents a formally verified Lightning Network.
% TODO: one sentence on the scope and the tool of [9] and how it differs from our work
% (keep it consistent with the corresponding sentence of the introduction).

Kiayias and Litos~\cite{kiayias2020} give a composable security treatment of the Lightning
Network: they define an ideal functionality for a payment network and prove, in the universal
composability framework, that the Lightning protocol realises it. Their guarantee is
computational and composable, and their proof is by hand. Ours are symbolic but machine-checked
by \tamarin; they include bypass resistance for routes of any length (for single-offer
payments), and they contrast the wormhole under HTLCs with its absence under PTLC-style locks on
a fixed route.

Malavolta et al.~\cite{malavolta2019} describe the wormhole attack on HTLC-based multi-hop
payments and propose anonymous multi-hop locks as a countermeasure; their instantiation with
adaptor signatures is the basis of PTLCs. Adaptor signatures also underlie the generalised
channels of Aumayr et al.~\cite{aumayr2021}, built from limited blockchain scripts. Our main
model reproduces the wormhole (\code{Wormhole_Fee_Theft_Reachable}, Sec.~\ref{sec:wormhole}),
and our auxiliary models check the countermeasure symbolically on a fixed route
(Sec.~\ref{sec:ptlc}), with a hash chain in place of the adaptor-signature algebra.

Smith et al.~\cite{smith2022} formalise agent-skipping attacks on message-forwarding protocols;
the wormhole is an instance of such an attack. Our bypass theorems (Sec.~\ref{sec:bypass}) state
the corresponding security property for Lightning routing: a forwarder at positions 1 to 6 of a
route, or at any position for single-offer payments, can be skipped only if some long-term key
is compromised.

Backes et al.~\cite{backes2017} reason about liveness in cryptographic protocols with a progress
semantics, in which honest processes that can continue without input from the adversary
eventually do, and apply it to fair exchange. We impose liveness as trace restrictions instead
(Sec.~\ref{par:liveness-scope}).

\section{Limitations}\label{sec:limitations}

\begin{itemize}
\item \emph{Abstract timing, liveness as restrictions.} The main model has no clock:
  \code{Open_Timeout} lets a channel time out at any point allowed by the restrictions, and T1, T2
  and T3 (the last two only for forwarders not compromised before their incoming HTLC times out)
  stand in for CLTV deadlines and timely claiming. As restrictions, they remove every trace that
  violates them, including traces in which an uncompromised forwarder could not have claimed in
  time; results proved under them say nothing about such runs (Sec.~\ref{par:liveness-scope}).
  Only the 24 lemmas of \texttt{make general}, among them Theorem~\ref{thm:anyhop}, are proved
  with all three removed; the necessity runs remove one at a time, and T3 is necessary for
  intermediary safety (Sec.~\ref{sec:timing}). Concrete block heights appear only in the auxiliary
  CLOCK and CLTV models, which are not composed with \code{lightning.spthy}. The progress
  semantics of Backes et al.~\cite{backes2017} would keep the traces in which an honest forwarder
  cannot act instead of pruning them. \tamarin{} 1.10 implements it for \sapic{} processes (option
  \code{translation-progress}), so the routing layer would have to be written as \sapic{}
  processes; we leave this to future work.
\item \emph{Single-offer premise.} The any-length theorem assumes that every HTLC for the
  payment hash is offered on the sender's channel; when a hash is offered more than once, only
  the fixed-position theorems (positions 1 to 6) apply, and these are proved under T1, T2 and T3.
  The conclusion of all bypass theorems says that some long-term key was compromised, not whose
  or when.
\item \emph{No end-to-end amounts.} Each \code{Forward_HTLC} reads its incoming amount, fee and
  outgoing amount from the network and requires the incoming amount to be the outgoing amount
  plus the fee, but the offer, the invoice and the \code{Route} fact carry no amount. Amounts are
  therefore not chained along the path, and we make no end-to-end claims about value, e.g., that
  the receiver obtains the invoiced amount or that the sender pays exactly the path fees.
\item \emph{HTLCs are not part of the commitment transactions.} The channel layer and the
  routing layer share the theory, the adversary and the channel topology (HTLCs are added only
  on opened channels, \code{HTLC_On_Opened_Channel}), but HTLC outputs are not folded into the
  commitment states, and settling or refunding an HTLC does not update the channel state. That
  each HTLC is resolved at most once is assumed as the restriction \code{OneResolutionPerHTLC},
  not derived from the channel layer; we do not claim a composition theorem. A channel carries
  at most one pending HTLC at a time (\code{Mutex(ptr)}).
\item \emph{Routing abstraction.} HTLC offers travel as the private linear fact \code{Route},
  not as network messages, so the network adversary cannot inject, replay or modify offers
  between peers; it chooses amounts and fees, controls the delivery of invoices and can
  compromise keys. The adversary never learns a preimage (\code{Preimage_Never_Known}), so in
  the main model the wormhole comes from the explicit rule \code{AttackerSkip}, which requires
  the parties at both ends of the skip to be compromised, rather than from adversary knowledge.
  (The auxiliary PTLC and HTLC models have no such rule.)
\item \emph{No onion routing, no privacy.} Every forwarder knows its predecessor, its successor
  and the payment hash, and the next hop is any party with which the forwarder shares a channel.
  We make no privacy claims.
\item \emph{PTLC model.} The countermeasure is checked only on the fixed route
  S$\,\to\,$H1$\,\to\,$H2$\,\to\,$H3$\,\to\,$R, with honest lock set-up, an honest sender and a
  hash chain in place of the adaptor-signature algebra (Sec.~\ref{sec:ptlc}).
\end{itemize}

\section{Conclusion}\label{sec:conclusion}

We presented one \tamarin{} theory of the Lightning Network that covers the channel layer and
multi-hop HTLC routing over routes of any length, against a Dolev-Yao adversary that can also
compromise long-term keys. We proved that only cheating or compromised parties are punished,
that settlements are authenticated by the receiver's invoice and that preimages stay secret;
under three liveness restrictions, an uncompromised intermediary never has its outgoing HTLC
redeemed while its incoming one is refunded, and payments are atomic in runs without key
compromise. The assumption that intermediaries claim in time is necessary: without it, \tamarin{}
finds a run without any compromise in which an honest intermediary loses. The wormhole is a
\tamarin-checked trace in which two compromised parties bypass an honest one, which is refunded
and earns no fee. Conversely, a bypass requires
some key compromise: under the liveness restrictions at positions 1 to 6 of any route, and, for
payments whose hash is offered on a single channel, at every position even without them. On a fixed route, per-hop locks in the style of PTLCs rule out the
bypass of an honest hop when the sender is honest, whereas a shared hash lock admits it. Next,
we want to replace the liveness restrictions by the progress semantics of~\cite{backes2017}, so
that runs in which an honest party cannot act are analysed rather than pruned, and to derive
\code{OneResolutionPerHTLC} from the channel layer instead of assuming it.

\section*{Artifact Availability}

The artifact is available as supplementary material.
% TODO: artifact URL for the camera-ready version (anonymised for review).
It consists of \code{lightning.spthy} (main model: channel layer, multi-hop routing and the
wormhole; 37 rules), \code{lightning_aux.spthy} (five auxiliary models, selected by
\code{-D=PTLC}, \code{-D=HTLC}, \code{-D=CLOCK}, \code{-D=CLTV} and \code{-D=TIMEOUT}) and a
\code{Makefile}. \texttt{make all} runs, in this order, \texttt{make verify}
(\code{lightning.spthy}, full model: 55/55 lemmas verified), \texttt{make general}
(Sec.~\ref{sec:bypass}: 24/24), \texttt{make aux} (\code{lightning_aux.spthy}: 21/21) and
\texttt{make necessity} (Sec.~\ref{sec:timing}); Table~\ref{tab:results} lists what each target
proves. All results were obtained with \tamarin{} 1.10 and reproduced on a laptop with a 6\,GB
heap (\texttt{MEM=6G}); \texttt{make verify-one} proves one lemma per run for machines with less
memory. \tamarin{} 1.10 reports the \maude{} version used, 3.2, as unsupported; all results above
were obtained with it. Every run also reports that \tamarin's derivation checks, part of its
well-formedness checks of the model, timed out; this warning does not concern the proofs.
